%%%%%%%%%%%%%%%%%%%%%%%%%%%%%%%%%%%%%%%%%%%%%%%%%%%%%%%%%%%%%%%%%%%%%%%%
%%%%%%%%%%%%%%%%%%%%%% Simple LaTeX CV Template %%%%%%%%%%%%%%%%%%%%%%%%
%%%%%%%%%%%%%%%%%%%%%%%%%%%%%%%%%%%%%%%%%%%%%%%%%%%%%%%%%%%%%%%%%%%%%%%%

%%%%%%%%%%%%%%%%%%%%%%%%%%%%%%%%%%%%%%%%%%%%%%%%%%%%%%%%%%%%%%%%%%%%%%%%
%% NOTE: If you find that it says                                     %%
%%                                                                    %%
%%                           1 of ??                                  %%
%%                                                                    %%
%% at the bottom of your first page, this means that the AUX file     %%
%% was not available when you ran LaTeX on this source. Simply RERUN  %%
%% LaTeX to get the ``??'' replaced with the number of the last page  %%
%% of the document. The AUX file will be generated on the first run   %%
%% of LaTeX and used on the second run to fill in all of the          %%
%% references.                                                        %%
%%%%%%%%%%%%%%%%%%%%%%%%%%%%%%%%%%%%%%%%%%%%%%%%%%%%%%%%%%%%%%%%%%%%%%%%

%%%%%%%%%%%%%%%%%%%%%%%%%%%% Document Setup %%%%%%%%%%%%%%%%%%%%%%%%%%%%

% Don't like 10pt? Try 11pt or 12pt
\documentclass[10pt]{article}

% The automated optical recognition software used to digitize resume
% information works best with fonts that do not have serifs. This
% command uses a sans serif font throughout. Uncomment both lines (or at
% least the second) to restore a Roman font (i.e., a font with serifs).
%\usepackage{times}
%\renewcommand{\familydefault}{\sfdefault}

% This is a helpful package that puts math inside length specifications
\usepackage{calc}
\usepackage{comment}
\usepackage{graphicx}

\usepackage{changepage}

% Simpler bibsection for CV sections
% (thanks to natbib for inspiration)
\makeatletter
\newlength{\bibhang}
\setlength{\bibhang}{1em} %1em}
\newlength{\bibsep}
 {\@listi \global\bibsep\itemsep \global\advance\bibsep by\parsep}
\newenvironment{bibsection}%
        {\begin{enumerate}{}{%
%        {\begin{list}{}{%
       \setlength{\leftmargin}{\bibhang}%
       \setlength{\itemindent}{-\leftmargin}%
       \setlength{\itemsep}{\bibsep}%
       \setlength{\parsep}{\z@}%
        \setlength{\partopsep}{0pt}%
        \setlength{\topsep}{0pt}}}
        {\end{enumerate}\vspace{-.6\baselineskip}}
%        {\end{list}\vspace{-.6\baselineskip}}
\makeatother

% Layout: Puts the section titles on left side of page
\reversemarginpar

%
%         PAPER SIZE, PAGE NUMBER, AND DOCUMENT LAYOUT NOTES:
%
% The next \usepackage line changes the layout for CV style section
% headings as marginal notes. It also sets up the paper size as either
% letter or A4. By default, letter was used. If A4 paper is desired,
% comment out the letterpaper lines and uncomment the a4paper lines.
%
% As you can see, the margin widths and section title widths can be
% easily adjusted.
%
% ALSO: Notice that the includefoot option can be commented OUT in order
% to put the PAGE NUMBER *IN* the bottom margin. This will make the
% effective text area larger.
%
% IF YOU WISH TO REMOVE THE ``of LASTPAGE'' next to each page number,
% see the note about the +LP and -LP lines below. Comment out the +LP
% and uncomment the -LP.
%
% IF YOU WISH TO REMOVE PAGE NUMBERS, be sure that the includefoot line
% is uncommented and ALSO uncomment the \pagestyle{empty} a few lines
% below.
%

%% Use these lines for letter-sized paper
\usepackage[paper=letterpaper,
            %includefoot, % Uncomment to put page number above margin
            marginparwidth=1.2in,     % Length of section titles
            marginparsep=.05in,       % Space between titles and text
            margin=.5in,               % 1 inch margins
            includemp]{geometry}


%% Use these lines for A4-sized paper
%\usepackage[paper=a4paper,
%            %includefoot, % Uncomment to put page number above margin
%            marginparwidth=30.5mm,    % Length of section titles
%            marginparsep=1.5mm,       % Space between titles and text
%            margin=25mm,              % 25mm margins
%            includemp]{geometry}

%% More layout: Get rid of indenting throughout entire document
\setlength{\parindent}{0in}

\usepackage[shortlabels]{enumitem}

%% Reference the last page in the page number
%
% NOTE: comment the +LP line and uncomment the -LP line to have page
%       numbers without the ``of ##'' last page reference)
%
% NOTE: uncomment the \pagestyle{empty} line to get rid of all page
%       numbers (make sure includefoot is commented out above)
%
\usepackage{fancyhdr,lastpage}
\pagestyle{fancy}
%\pagestyle{empty}      % Uncomment this to get rid of page numbers
\fancyhf{}\renewcommand{\headrulewidth}{0pt}
\fancyfootoffset{\marginparsep+\marginparwidth}
\newlength{\footpageshift}
\setlength{\footpageshift}
          {0.5\textwidth+0.5\marginparsep+0.5\marginparwidth-2in}
% \lfoot{\hspace{\footpageshift}%
%        \parbox{4in}{\, \hfill %
%                     \arabic{page} of \protect\pageref*{LastPage} % +LP
% %                    \arabic{page}                               % -LP
%                     \hfill \,}}

% Finally, give us PDF bookmarks
\usepackage{color,hyperref}
\definecolor{darkblue}{rgb}{0.0,0.0,0.3}
\hypersetup{colorlinks,breaklinks,
            linkcolor=darkblue,urlcolor=darkblue,
            anchorcolor=darkblue,citecolor=darkblue}

%%%%%%%%%%%%%%%%%%%%%%%% End Document Setup %%%%%%%%%%%%%%%%%%%%%%%%%%%%


%%%%%%%%%%%%%%%%%%%%%%%%%%% Helper Commands %%%%%%%%%%%%%%%%%%%%%%%%%%%%

% The title (name) with a horizontal rule under it
% (optional argument typesets an object right-justified across from name
%  as well)
%
% Usage: \makeheading{name}
%        OR
%        \makeheading[right_object]{name}
%
% Place at top of document. It should be the first thing.
% If ``right_object'' is provided in the square-braced optional
% argument, it will be right justified on the same line as ``name'' at
% the top of the CV. For example:
%
%       \makeheading[\emph{Curriculum vitae}]{Your Name}
%
% will put an emphasized ``Curriculum vitae'' at the top of the document
% as a title. Likewise, a picture could be included:
%
%   \makeheading[\includegraphics[height=1.5in]{my_picutre}]{Your Name}
%
% the picture will be flush right across from the name.
\newcommand{\makeheading}[2][]%
        {\hspace*{-\marginparsep minus \marginparwidth}%
         \begin{minipage}[t]{\textwidth+\marginparwidth+\marginparsep}%
             {\large \bfseries #2 \hfill #1}\\[-0.15\baselineskip]%
                 \rule{\columnwidth}{1pt}%
         \end{minipage}}

% The section headings
%
% Usage: \section{section name}
\renewcommand{\section}[1]{\pagebreak[3]%
    \hyphenpenalty=10000%
    \vspace{1.3\baselineskip}%
    \phantomsection\addcontentsline{toc}{section}{#1}%
    \noindent\llap{\scshape\smash{\parbox[t]{\marginparwidth}{\raggedright #1}}}%
    \vspace{-\baselineskip}\par}

% An itemize-style list with lots of space between items
\newenvironment{outerlist}[1][\enskip\textbullet]%
        {\begin{itemize}[#1,leftmargin=*]}{\end{itemize}%
         \vspace{-.6\baselineskip}}

% An environment IDENTICAL to outerlist that has better pre-list spacing
% when used as the first thing in a \section
\newenvironment{lonelist}[1][\enskip\textbullet]%
        {\begin{list}{#1}{%
        \setlength{\partopsep}{0pt}%
        \setlength{\topsep}{0pt}}}
        {\end{list}\vspace{-.6\baselineskip}}

% An itemize-style list with little space between items
\newenvironment{innerlist}[1][\enskip\textbullet]%
        {\begin{itemize}[#1,leftmargin=*,parsep=0pt,itemsep=0pt,topsep=0pt,partopsep=0pt]}
        {\end{itemize}}

% An environment IDENTICAL to innerlist that has better pre-list spacing
% when used as the first thing in a \section
\newenvironment{loneinnerlist}[1][\enskip\textbullet]%
        {\begin{itemize}[#1,leftmargin=*,parsep=0pt,itemsep=0pt,topsep=0pt,partopsep=0pt]}
        {\end{itemize}\vspace{-.6\baselineskip}}

\newenvironment{myinnerlist}[1][\enskip\textbullet]%
        {\begin{itemize}[#1,leftmargin=*,parsep=0pt,itemsep=3pt,topsep=0pt,partopsep=0pt]}
        {\end{itemize}\vspace{.6\baselineskip}}

% To add some paragraph space between lines.
% This also tells LaTeX to preferably break a page on one of these gaps
% if there is a needed pagebreak nearby.
\newcommand{\blankline}{\quad\pagebreak[3]}
\newcommand{\halfblankline}{\quad\vspace{-0.5\baselineskip}\pagebreak[3]}

% Uses hyperref to link DOI
\newcommand\doilink[1]{\href{http://dx.doi.org/#1}{#1}}
\newcommand\doi[1]{doi:\doilink{#1}}

% For \url{SOME_URL}, links SOME_URL to the url SOME_URL
\providecommand*\url[1]{\href{#1}{#1}}
% Same as above, but pretty-prints SOME_URL in teletype fixed-width font
\renewcommand*\url[1]{\href{#1}{\texttt{#1}}}

% For \email{ADDRESS}, links ADDRESS to the url mailto:ADDRESS
\providecommand*\email[1]{\href{mailto:#1}{#1}}
% Same as above, but pretty-prints ADDRESS in teletype fixed-width font
%\renewcommand*\email[1]{\href{mailto:#1}{\texttt{#1}}}

%\providecommand\BibTeX{{\rm B\kern-.05em{\sc i\kern-.025em b}\kern-.08em
%    T\kern-.1667em\lower.7ex\hbox{E}\kern-.125emX}}
%\providecommand\BibTeX{{\rm B\kern-.05em{\sc i\kern-.025em b}\kern-.08em
%    \TeX}}
\providecommand\BibTeX{{B\kern-.05em{\sc i\kern-.025em b}\kern-.08em
    \TeX}}
\providecommand\Matlab{\textsc{Matlab}}


\usepackage{xcolor}
\newcommand{\localtextbulletone}{\textcolor{gray}{\raisebox{.45ex}{\rule{.8ex}{.8ex}}}}

%%%%%%%%%%%%%%%%%%%%%%%% End Helper Commands %%%%%%%%%%%%%%%%%%%%%%%%%%%

%%%%%%%%%%%%%%%%%%%%%%%%% Begin CV Document %%%%%%%%%%%%%%%%%%%%%%%%%%%%

\begin{document}
\makeheading{{\Huge Siddhartha}\hfill{\large FPGA/Software Engineer }}

\section{Contact Information}

% NOTE: Mind where the & separators and \\ breaks are in the following
%       table.
%
% ALSO: \rcollength is the width of the right column of the table
%       (adjust it to your liking; default is 1.85in).
%
\newlength{\rcollength}\setlength{\rcollength}{1.4in}%
%
\begin{tabular}[t]{@{}p{\textwidth-\rcollength-1.8in}p{2\rcollength}}
%\href{http://www.cse.osu.edu/}%
%     {Department of Computer Science and Engineering} & \\
%\href{http://www.osu.edu/}{The Ohio State University}
\includegraphics[width=12px]{email_logo}\ \ \ \textcolor{black}{\textsf{sidmontu@gmail.com}} & 
	\includegraphics[width=12px]{telephone_logo}\ \ \ {\textsf{+65 9857-4171}} \\
    \includegraphics[width=12px]{globe_logo}\ \ \ \textcolor{black}{\href{https://sidmontu.github.io}{\textsf{https://sidmontu.github.io}}} & 
	\includegraphics[width=12px]{github_logo}\ \ \ \raisebox{0.3\height}{\textcolor{black}{\href{https://github.com/sidmontu}{\textsf{github.com/sidmontu}}}} \\
	\multicolumn{2}{@{}}{\includegraphics[width=12px]{home_logo}\ \ \ {\textsf{7, Woodsvale Condominium, \#06-13, Woodlands Drive 72, Singapore 738092}}} \\
\end{tabular}

\section{Interests}

FPGAs, Machine Learning Engineering, HPC, RF Communications

\section{Professional Summary}

Experienced mid-career software engineer with a FPGA-focused research
background looking for a new challenge. Especially interested in projects at
the crossroads of hardware acceleration (FPGAs/ASICs) and modern deep learning
techniques. Referees available on request.
% background looking for a fresh challenge. Especially interested in applied AI
% and ML infrastructure projects (e.g. MLOps), with/without hardware
% acceleration (e.g. FPGAs/ASICs) as an auxiliary goal. Referees on request.

{
    \hspace*{-\marginparsep minus \marginparwidth}%
    \begin{minipage}[t]{\textwidth+\marginparwidth+\marginparsep}%
         \rule{\columnwidth}{1pt}\\[-1.4\baselineskip]%
         \begin{center}{\large \bfseries Work Experience}\end{center}
         \vspace{-0.18in}
         \rule{\columnwidth}{1pt}\\[-1.4\baselineskip]%
    \end{minipage}
}\\[.5\baselineskip]


{
    \hspace*{-\marginparsep minus \marginparwidth}%
    \begin{minipage}[t]{\textwidth+\marginparwidth+\marginparsep}%

        \textbf{Machine Learning Engineer (Contract), SAP Singapore} \hfill {Jun 2021 to Present}

        \vspace{0.1in}
        \begin{myinnerlist}
            \renewcommand{\labelitemi}{\localtextbulletone}

            \item Building enterprise search solutions for internal stakeholders using Elasticsearch/Opensearch stacks.
            \item Implementing DNN-based semantic-similarity algorithms for information retrieval and ranking of results.
            \item MLOps tooling: setting up labeling, data cleaning, and evaluation pipelines for efficient data-centric ML implementations.
            \item Tools/frameworks used (non-exhaustive): [Elastic$\vert$Open]search, Docker, FastAPI, Kubernetes, Poetry, Kibana, etc.

        \end{myinnerlist}

        \textbf{CTO/Founder, Inpact Technologies Pte Ltd (inPact.ai), Singapore} \hfill {May 2020 to Sep 2021}

        \vspace{0.1in}
        \begin{myinnerlist}
            \renewcommand{\labelitemi}{\localtextbulletone}

            \item AS CTO, planned and executed in-house product development
                across all aspects of machine learning, frontend, backend, and
                deployment. Hired freelancers and interns to achieve two
                product release milestones within 7 months.

            \item Frontend using \href{https://react-redux.js.org/}{\bf
                React-Redux} + \href{https://sass-lang.com/}{\bf SASS} stack;
                backend using \href{https://aws.amazon.com/serverless/}{\bf
                serverless framework}; deployment on AWS (EC2, DynamoDB,
                Cognito, S3, etc); ML models trained using PyTorch and
                \href{https://spacy.io/}{\bf spaCy}, REST API deployment through
                \href{https://www.cortex.dev/}{\bf cortex}.
                % Hands-on experience with technologies/frameworks such as
                % ReactJS/Redux/SASS, serverless, spaCy, and more.

            \item Trained/fine-tuned DNNs (e.g. residual CNNs, BERT) for NLP
                tasks like text classification and named entity recognition.

            % \item Handled business functions such as customer acquisition,
            %     investor relations, fundraising,
            %     marketing, and more.
            \item Raised S\$75K from Entrepreneur First, an international VC
                firm funded by some of the top investors in the world.

        \end{myinnerlist}

        \textbf{Posdoctoral Research Associate, University of Sydney, Australia} \hfill {Oct 2017 to Dec 2019}

        \vspace{0.1in}
        \begin{myinnerlist}
            \renewcommand{\labelitemi}{\localtextbulletone}

            \item Research \& development on a ``High-Speed Machine Learning
                for RF Communications'' project commissioned by the Australian
                Defense Agency, leading to a peer-reviewed publication at
                \href{https://ieeexplore.ieee.org/document/8599833}{\bf MILCOM
                2018}. PI:
                \href{https://www.sydney.edu.au/engineering/about/our-people/academic-staff/philip-leong.html}{\bf
                Prof. Philip Leong}.
            \item Contribution to other FPGA-based research projects in the lab on
                topics such as
                \href{https://ieeexplore.ieee.org/document/8742277}{\bf on-chip
                training},
                \href{https://dl.acm.org/doi/abs/10.1145/3373087.3375303}{\bf logic
                cell architecture design}, and
                \href{https://ieeexplore.ieee.org/abstract/document/9150443}{\bf automatic
                modulation classification}.
                % Total 5 research papers/posters published at peer-reviewed conferences.
            \item Extensive hands-on experience with technologies such as
                Vivado[HLS], RFSoC, PYNQ, Tensorflow, etc during stint at the
                lab. Undertook sysadmin duties to manage lab resources, and
                advised students on their final-year research projects.
            % \item Supervision \& guidance to undergraduate students on their
            %     final-year research projects.
            % \item Sysadmin duties on available compute resources in the lab.

        \end{myinnerlist}
        \vspace{-0.15in}
    \end{minipage}
}\\[.15\baselineskip]

{
    \hspace*{-\marginparsep minus \marginparwidth}%
    \begin{minipage}[t]{\textwidth+\marginparwidth+\marginparsep}%
         \rule{\columnwidth}{1pt}\\[-1.4\baselineskip]%
         \begin{center}{\large \bfseries Education \& Certifications}\end{center}
         \vspace{-0.18in}
         \rule{\columnwidth}{1pt}\\[-1.4\baselineskip]%
     \end{minipage}}\\[.5\baselineskip]

{
    \hspace*{-\marginparsep minus \marginparwidth}%
    \begin{minipage}[t]{\textwidth+\marginparwidth+\marginparsep}%

        {\textbf{Nanyang Technological University (NTU), Singapore}} \hfill {Jan 2013 to Feb 2019}

        \vspace{0.1in}
        \begin{innerlist}
            \item[] PhD, {Computer Science \& Engineering},
                Advisor: \href{https://nachiket.github.io/}{\bf Prof. Nachiket Kapre}
            \item[] Dissertation: \href{https://dr.ntu.edu.sg/handle/10220/47803}{\bf Dataflow Optimized Overlays for FPGAs}
            \item[]
                \begin{myinnerlist}
                \renewcommand{\labelitemi}{\localtextbulletone}
                \item Developed a 16x16 mesh token dataflow overlay
                    architecture finely-tuned for efficient mapping on Arria 10
                    FPGAs.
                \item Exploited statically-extracted instruction-criticality
                    information to enable out-of-order execution inside each
                    processing element; improved performance by up to
                    2.4$\times$ over existing in-order processor architectures.
                \item Verilog for hardware implementation; verilator for
                    testbenches and cycle-accurate simulations; C++ compiler to
                    optimize and map dataflow graphs to instructions inside
                    each PE; python/R/shell scripting for data wrangling tasks.
            \end{myinnerlist}
            \vspace{-0.06in}
            \item[] Teaching Assistant: Programmable System on Chip (CE4054), Embedded Software Development (CE4052)
        \end{innerlist}

        \vspace{.1in}
        {\textbf{Imperial College London, United Kingdom}} \hfill {Oct 2009 to Jun 2012}
        \begin{outerlist}
        \item[] Bachelors of Engineering (BEng), {Electrical \& Electronics Engineering}
        \end{outerlist}

        \vspace{.1in}
        {\bf \href{https://www.coursera.org/specializations/machine-learning-engineering-for-production-mlops}{Machine Learning Engineering for Production Specialization, DeepLearning.AI, Coursera}}\hfill{Ongoing}\\
        % \begin{innerlist}
        %     \item[] 4-module specialization offered through the Coursera MOOC platform.
        %         Focuses on best in-industry practices for deploying data-centric machine
        %         learning systems. Instructors include Andrew Ng (founder DeepLearning.AI
        %         and Coursera), Robert Crowe (Tensorflow Developer Engineer, Google),
        %         Laurence Moroney (Lead AI Advocate, Google), and more.
        % \end{innerlist}

        \vspace{-.1in}
        {\bf \href{https://www.coursera.org/account/accomplishments/specialization/YP6NMPR4X3M9}{Data Science Specialization, John Hopkins University, Coursera}}\hfill{Apr 2016}\\[-0.1in]
        % {\color{darkgray}\small Offered through the Coursera MOOC platform}
        % \begin{innerlist}
        %     \item[] 10-module specialization that covers concepts and tools essential for
        %         building effective data science pipelines. Capstone project on building a
        %         text-prediction R web-app, similar to the predictive text technology found
        %         in mobile phones today.
        % \end{innerlist}

    \end{minipage}
}\\[.15\baselineskip]

\newpage

%%%%%%%%%%%%%%%%%%%%%%%%%%%%%%%%%%%%%%%%%%%%%%%%%%%%%%%%%%%%%%%%%%%%%%%%%%%%%%%%%%%%%%%%%%%%%%%
%%%%%%%%%%%%%%%%%%%%%%%%%%%%%%%%%%%%%%%% Publications %%%%%%%%%%%%%%%%%%%%%%%%%%%%%%%%%%%%%%%%%
%%%%%%%%%%%%%%%%%%%%%%%%%%%%%%%%%%%%%%%%%%%%%%%%%%%%%%%%%%%%%%%%%%%%%%%%%%%%%%%%%%%%%%%%%%%%%%%

{
    \hspace*{-\marginparsep minus \marginparwidth}%
    \begin{adjustwidth}{-1.25in}{}
    % \begin{minipage}[t]{\textwidth+\marginparwidth+\marginparsep}%

        % \section{Peer-Reviewed Publications}
        {\Large Peer-Reviewed Publications}\\
        {
            \hspace*{-\marginparsep}%
            \begin{minipage}[t]{1.25in+\textwidth+\marginparsep}%
                \rule{\columnwidth}{1pt}\\[-1.4\baselineskip]%
            \end{minipage}
        }\\[.5\baselineskip]

        \textbf{\large Full Papers (Conferences/Journals)}
        \vspace{0in}
        \begin{enumerate}
            \item Seyedramin Rasoulinezhad, {\bf Siddhartha}, Hao Zhou, Lingli Wang, David Boland, and Philip Leong ``LUXOR: An FPGA Logic Cell Architecture for Efficient Compressor Tree Implementations" \emph{Proceedings of the 2020 ACM/SIGDA International Symposium on Field-Programmable Gate Arrays (FPGA)}, February 2020\newline[\href{https://doi.org/10.1145/3373087.3375303}{\texttt{DOI: 10.1145/3373087.3375303}}]
            \item {\bf Siddhartha}, and Nachiket Kapre ``DaCO: A High-Performance Token Dataflow Coprocessor Overlay for FPGAs" \emph{International Conference on Field-Programmable Technology}, December 2018\newline[\href{https://doi.org/10.1109/FPT.2018.00032}{\texttt{DOI: 10.1109/FPT.2018.00032}}]
            \item {\bf Siddhartha}, Yee Hui Lee, Duncan Moss, Julian Faraone, Perry Blackmore, Daniel Salmond, David Boland, and Philip Leong ``Long Short-Term Memory for Radio Frequency Spectral Prediction and its Real-Time FPGA Implementation" \emph{IEEE Military Communications Conference (MILCOM)}, October 2018\newline[\href{https://doi.org/10.1109/MILCOM.2018.8599833}{\texttt{DOI: 10.1109/MILCOM.2018.8599833}}]
            \item {\bf Siddhartha}, Nachiket Kapre ``Hoplite-Q: Priority-Aware Routing in FPGA Overlay NoCs" \emph{IEEE 26th Annual International Symposium on Field-Programmable Custom Computing Machines}, May 2018\newline[\href{https://doi.org/10.1109/FCCM.2018.00012}{\texttt{DOI: 10.1109/FCCM.2018.00012}}]
            \item Gopalakrishna Hegde, {\bf Siddhartha}, Nachiket Kapre ``CaffePresso: Accelerating Convolutional Networks on Embedded SoCs'' \emph{ACM Transactions on Embedded Computing Systems (TECS)}, January 2018\newline[\href{https://doi.org/10.1145/3105925}{\texttt{DOI: 10.1145/3105925}}]
            \item {\bf Siddhartha}, Nachiket Kapre ``eBSP: Managing NoC traffic for BSP workloads on the 16-core Adapteva Epiphany-III Processor." \emph{Design, Automation, and Test in Europe}, March 2017\newline[\href{https://doi.org/10.23919/DATE.2017.7926961}{\texttt{DOI: 10.23919/DATE.2017.7926961}}]
            \item Gopalakrishna Hegde, \textbf{Siddhartha}, Nachiappan Ramasamy, Nachiket Kapre\newline``CaffePresso: An Optimized Library for Deep Learning on Embedded Accelerator-based platforms." \emph{International Conference on Compilers, Architecture, and Synthesis for Embedded Systems}, October 2016 (Best Paper Award)\newline[\href{https://doi.org/10.1145/2968455.2968511}{\texttt{DOI: 10.1145/2968455.2968511}}]
            \item Pradeep Moorthy, {\bf Siddhartha}, and Nachiket Kapre ``A Case for Embedded FPGA-based SoCs for Energy-Efficient Acceleration of Graph Problems."\newline\emph{Supercomputing Frontiers 2015}, March 2015\newline[\href{http://dx.doi.org/10.14529/jsfi150307}{\texttt{DOI: 10.14529/jsfi150307}}]
        \end{enumerate}

        \vspace{0.1in}
        \textbf{\large Short Papers/Posters/Workshops}
        \begin{enumerate}
            \item Stephen Tridgell, David Boland, Philip Leong, Ryan Kastner, Alireza Khodamoradi, and {\bf Siddhartha} ``Real-time Automatic Modulation Classification using RFSoC" \emph{27th Reconfigurable Architectures Workshop (co-located with IPDPS 2020)}, May 2020\newline[\href{https://doi.org/10.1109/IPDPSW50202.2020.00021}{\texttt{DOI: 10.1109/IPDPSW50202.2020.00021}}]
            \item Stephen Tridgell, David Boland, Philip Leong, and {\bf Siddhartha} ``Real-time Automatic Modulation Classification" \emph{International Conference on Field-Programmable Technology}, December 2019 (Poster)\newline[\href{https://doi.org/10.1109/ICFPT47387.2019.00052}{\texttt{DOI: 10.1109/ICFPT47387.2019.00052}}]
            \item {\bf Siddhartha}, David Boland, Steve Wilton, Barry Flower, Perry Blackmore, and Philip Leong ``Simultaneous Inference and Training using On-FPGA Weight Perturbation Techniques" \emph{International Conference on Field-Programmable Technology}, December 2018 (Poster)\newline[\href{https://doi.org/10.1109/FPT.2018.00060}{\texttt{DOI: 10.1109/FPT.2018.00060}}]
            \item {\bf Siddhartha}, Nachiket Kapre ``Out-of-Order Dataflow Scheduling for FPGA Overlays." \emph{Overlay Architectures for FPGAs Workshop (co-located with FPGA 2017)}, February 2017 (Position Paper)\newline[\href{https://arxiv.org/abs/1704.08802}{\texttt{DOI: arXiv:1705.02734}}]
            \item Sidharth Maheshwari, Gourav Modi, {\bf Siddhartha}, Nachiket Kapre ``Vector FPGA Acceleration of 1-D DWT Computations using Sparse Matrix Skeletons." \emph{26th IEEE International Conference on Field-Programmable Logic and Applications}, August 2016 (Poster)\newline[\href{https://doi.org/10.1109/FPL.2016.7577361}{\texttt{DOI: 10.1109/FPL.2016.7577361}}]
            \item {\bf Siddhartha}, Nachiket Kapre ``Communication Optimization for the 16-core Epiphany Floating-Point Processor Array." \emph{24th IEEE International Symposium on Field-Programmable Custom Computing Machines}, May 2016 (Short Paper)\newline[\href{https://doi.org/10.1109/FCCM.2016.15}{\texttt{DOI: 10.1109/FCCM.2016.15}}]
            \item Gopalakrishna Hegde, {\bf Siddhartha}, Nachiappan Ramasamy, Vamsi Buddha, Nachiket Kapre ``Evaluating Embedded FPGA Accelerators for Deep Learning Applications." \emph{24th IEEE International Symposium on Field-Programmable Custom Computing Machines}, May 2016 (Short Paper)\newline[\href{https://doi.org/10.1109/FCCM.2016.14}{\texttt{DOI: 10.1109/FCCM.2016.14}}]
            \item Nachiket Kapre, Han Jianglei, Andrew Bean, Pradeep Moorthy, and {\bf Siddhartha} ``GraphMMU: Memory Management Unit for Sparse Graph Accelerators." \emph{22nd Reconfigurable Architectures Workshop (co-located with IPDPS)}, May 2015\newline[\href{https://doi.org/10.1109/IPDPSW.2015.101}{\texttt{DOI: 10.1109/IPDPSW.2015.101}}]
            \item {\bf Siddhartha}, Nachiket Kapre ``FPGA Acceleration of Irregular Iterative Computations using Criticality-Aware Dataflow Optimizations." \emph{International Symposium on Field-Programmable Gate Arrays}, February 2015 (Short Paper)\newline[\href{https://doi.org/10.1145/2684746.2689110}{\texttt{DOI: 10.1145/2684746.2689110}}]
            \item {\bf Siddhartha}, Nachiket Kapre ``Fanout Decomposition Dataflow Optimizations for FPGA-based Sparse LU Factorization." \emph{International Conference on Field-Programmable Technology}, December 2014 (Short Paper)\newline[\href{https://doi.org/10.1109/FPT.2014.7082787}{\texttt{DOI: 10.1109/FPT.2014.7082787}}]
            \item {\bf Siddhartha}, Nachiket Kapre ``Heterogeneous Dataflow Architectures for FPGA-based Sparse LU Factorization." \emph{The International Conference on Field Programmable Logic and Applications}, September 2014 (Short Paper)\newline[\href{https://doi.org/10.1109/FPL.2014.6927401}{\texttt{DOI: 10.1109/FPL.2014.6927401}}]
            \item Nachiket Kapre, {\bf Siddhartha} ``Limits of Statically-Scheduled Token Dataflow Processing." \emph{International workshop on Data-Flow Models (DFM) for Extreme Scale Computing (co-located with PACT 2014)}, August 2014\newline[\href{https://doi.org/10.1109/DFM.2014.21}{\texttt{DOI: 10.1109/DFM.2014.21}}]
            \item {\bf Siddhartha}, Nachiket Kapre ``Breaking Sequential Dependencies in FPGA-based Sparse LU Factorization." \emph{International Symposium on Field Programmable Custom Computing Machines}, May 2014 (Short Paper)\newline[\href{https://doi.org/10.1109/FCCM.2014.26}{\texttt{DOI: 10.1109/FCCM.2014.26}}]
        \end{enumerate}

    % \end{minipage}
    \end{adjustwidth}
}% \\[.15\baselineskip]

\end{document}

%%%%%%%%%%%%%%%%%%%%%%%%%% End CV Document %%%%%%%%%%%%%%%%%%%%%%%%%%%%%

%----------------------------------------------------------------------%
% The following is copyright and licensing information for
% redistribution of this LaTeX source code; it also includes a liability
% statement. If this source code is not being redistributed to others,
% it may be omitted. It has no effect on the function of the above code.
%----------------------------------------------------------------------%
% Copyright (c) 2007, 2008, 2009, 2010, 2011 by Theodore P. Pavlic
%
% Unless otherwise expressly stated, this work is licensed under the
% Creative Commons Attribution-Noncommercial 3.0 United States License. To
% view a copy of this license, visit
% http://creativecommons.org/licenses/by-nc/3.0/us/ or send a letter to
% Creative Commons, 171 Second Street, Suite 300, San Francisco,
% California, 94105, USA.
%
% THE SOFTWARE IS PROVIDED "AS IS", WITHOUT WARRANTY OF ANY KIND, EXPRESS
% OR IMPLIED, INCLUDING BUT NOT LIMITED TO THE WARRANTIES OF
% MERCHANTABILITY, FITNESS FOR A PARTICULAR PURPOSE AND NONINFRINGEMENT.
% IN NO EVENT SHALL THE AUTHORS OR COPYRIGHT HOLDERS BE LIABLE FOR ANY
% CLAIM, DAMAGES OR OTHER LIABILITY, WHETHER IN AN ACTION OF CONTRACT,
% TORT OR OTHERWISE, ARISING FROM, OUT OF OR IN CONNECTION WITH THE
% SOFTWARE OR THE USE OR OTHER DEALINGS IN THE SOFTWARE.
%----------------------------------------------------------------------%
